\setchapterpreamble[u]{\margintoc}
\chapter{Measurement of the Transfer Function}
\labch{tf-measurement}

If we have an optical system and want to model the way light passes through that system, either in order to enhance our
    estimate of the structure (source) that causes the output signal or to establish characteristics of that system, we 
    need to determine one of the transfer functions introduced in \refch{transfer-functions}. % TODO: reference <14-11-23> 
Of the transfer functions, it is most intuitive to directly measure the \ac{PSF} or more specifically the
    magnitude of the \ac{PSF}.
We disregard the phase change that the optical system may introduce, because in most set-ups it is practically 
    impossible to directly measure the imaginary part of the \ac{PSF}\sidenote[][*+12]{%
    It is possible in some cases to determine the imaginary part (i.e. the phase change \ac{PSF}) from the magnitude 
        \ac{PSF} algorithmically if we assume certain properties of the complex \ac{PSF}. 
    This algorithmic restoration of the phase is called phase-retrieval. However, in order for the phase to be 
        unambiguous, thus for the phase-retrieval algorithm to be stable, it is necessary that the phase retrieved 
        function to be non-symmetric with respect to the origin, which is often the case for the \ac{PSF}
        \sidecite[*+2]{goodman2005b}.
    There are other algorithms that for phase retrieval that leverage the fact that the prior knowledge of the \ac{PSF}
    \sidecite[*+3]{hanser2003}.
    } , since 
    the light sensors (In our case the \ac{CMOS}) only measure the energy that hits the sensor elements which is 
    equivalent for all possible phases of the \ac{EM} wave, since it only depends on the intensity of the \ac{EM} 
    wave~(\refchshort{light}). % TODO: Add closer reference. <14-11-23> 

For estimating the transfer function, it is necessary to make an acquisition of a known source signal.
It is possible to categorize the following methods of estimation into two parts.
At first, we will avoid creating a model of the whole source image \(\source\) and only leverage the knowledge of the
    objects that are being observed.
For this, a more comprehensive formulation of the image formation will be discussed in \refsec{image-formation-model} 
    and then the model-free estimate of the \ac{PSF} will be performed~(\refsecshort{psf-from-beads}).
Next, a complete model \(\model\) of the source signal \(\source\) will be determined~(\refsecshort{source-model}) and a 
    method for estimation of the \ac{OTF} using this model will be demonstrated.
Afterwards, one more estimate of the \ac{PSF} that leverages the model and recognizes the fact that the data is from a 
    \ac{SIM} acquisition will be performed~(\refsecshort{psf-from-sim})
Finally, we will return to the topic from \refch{sim-reconstruction} and compare the reconstructions that use the 
    transfer function estimates from this chapter~(\refsecshort{revised-reconstruction}).

\section{Image Formation Model}%
\labsec{image-formation-model}

% TODO: First step, noise. Measure noise from the background of the beads image. <14-11-23> 
Devising an image formation model means to formulate the manner in which, we expect light from the source to manifest 
    itself in the resulting sensed signal.
For reasoning about the later estimates of the transfer functions, it is necessary to express this image formation model
    mathematically.
It is sensible to proceed in the same chronological order as the light passing through the apparatus.
With this in mind, we will start off with the illumination of the focal plane.

\subsection{Illumination}%
\labsubsec{illumination}
The transfer function of the optical system, is given by measuring the way that light passes though the apparatus given 
    a uniformly illuminated focal plane.
However, we do not have data directly observed using a regular uniformly illuminated wide field microscopy system.
% TODO: Reference specific section <22-11-23> 
Our images are acquired with the harmonic illumination~(\refchshort{SR-microscopy}). 
This means that for a still\sidenote{Meaning that the source of the light is stationary in contrast to \emph{in 
    vivo} biological acquisitions, where this is not the case.} image, we have 9 wide field images with 3 different 
    pattern orientations and 3 different phases\sidenote[][*+5]{In accordance with the method that is used to generate 
    the patterns in the case of our \ac{SIM} setup, we are assuming that the phases differ by \(2\pi/3\) from one 
    another and do not determine them independently~(\refsecshort{IP-parameter-estimation}).} for each of those.

In order to account for this discrepancy, we will emulate the uniform illumination by summing all the images acquired
    with different illuminations together (this will also reduce noise~(\refsubsecshort{noise})).
Due to additivity of the optical system, this results in the same signal, as if an acquisition was made with the sum
    of the illumination patterns.
    For three patterns in one orientation \(\ip_-\), \(\ip_0\) and \(\ip_+\), where the phases are \(\phi_0
    - \tfrac{2\pi}{3}\), \(\phi_0\)
    and \(\phi_0 + \tfrac{2\pi}{3}\), let us, without loss of generality, assume that \(\phi_0 \equiv 0\) (since we can move the 
    origin as we please), simplifying the sum \(\ip_\Sigma\)~(\reffigshort{sum-of-ips}) along the line of the wave 
    vector of the harmonic pattern i.e. \(\vec{r} = \alpha(\vec{k}_\ip)^{\transpose}\), gives % TODO: Graphic of the line in the pattern <22-11-23> 
    \begin{equation*}
        \begin{split}
            \ip_\Sigma (\vec{r}) &= \sum_{i=1}^3 1 - \frac{m_i}{2} \cos\ab(2\pi \vec{k}_\ip \cdot \vec{r} 
                + \frac{2\pi i}{3}) \\
                               &= 3 - \sum_{i=1}^3 \frac{m}{2} \cos\ab(2\pi \vec{k}_\ip \cdot \vec{r} + 
                                \frac{2\pi i}{3}) \\
                               &\begin{split}
                                   = 3 &-\frac{m}{2} \cos\ab(2\pi \vec{k}_\ip \cdot \vec{r})
                                    \underbrace{\sum_{i=1}^3\ab(\cos\frac{2\pi i}{3})}_{=0} \\
                                       &- \frac{m}{2} \sin\ab(2\pi  \vec{k}_\ip \cdot \vec{r})
                                    \underbrace{\sum_{i=1}^3\ab(\sin\frac{2\pi i}{3})}_{=0},
                                \end{split}
        \end{split}
    \end{equation*}
    where, in the second step, we made the assumption that \(m_1 = m_2 = m_3\)\sidenote{This assumption holds 
    precisely if the illumination pattern is generated by a diffraction grating, which is our case, and also 
    theoretically if laser interference is used.}.

\begin{marginfigure}[*-16]
\marginpdftex{sum-of-patterns}
\caption{Sum of the illumination pattern intensity along the wave vector for the first orientation, results in a uniform 
    intensity.}
\labfig{sum-of-ips}
\end{marginfigure}

% TODO: Show that along another line, this only projects the wave vector on the line, and the phases do not change so
% the sum is also uniform <22-11-23> 

% TODO: Section about the images <16-11-23> 
% TODO: Figure: Add reconstructed (HR) image figure and LR image figures <16-11-23> 

% TODO: Second step, illumination... The sum <16-11-23> 
    % TODO: Side figure of the z-x axis cut of the illumination intensity cut <16-11-23> 
    % TODO: Note about additivity and why the 3× illuminated signal is the same as if we would have illuminated it all at
    % once <16-11-23> 

\subsection{Noise}%
\labsubsec{noise}

As a first step towards developing a faithful image formation model of a uniformly illuminated source, we need to 
    consider noise.

The mode of noise, that we will be particularly interested in is \emph{additive noise}, for which the strength of the 
    noise is added to the noise free source signal\sidenote[][*-4]{The location in the apparatus in which noise corruption 
        happens is particularly interesting for the purpose of 
        estimating the transfer function, because if the noise is added to the signal before the pass through the 
        optical system, it is degraded with the transfer function along with the signal and if the corruption 
        happens after, the noise affects the already degraded signal and is independent of the transfer function. 
    As a simplification measure, and in compliance with the major sources of noise, we will only regard noise that 
        corrupts the signal after the pass through the optical system.
        }:
\begin{equation} \labeq{additive-noise-model}
    \sensed(\vec{r}) = \overbrace{\noisefree(\vec{r})}^{\mathclap{\text{signal}}}
    + \underbrace{\noise(\vec{r})}_{\mathclap{\text{noise}}}.
\end{equation}
Additionally, in~\refeq{additive-noise-model} the noise \(\noise(\vec{r})\) has a specified distribution\sidenote[][*+6]{It is not 
    uncommon to choose to model the noise by a Gaussian distribution (\(\distribution{P} \equiv \distribution{N}\)) even
    if the Gaussian distribution is physically unrealistic for the given noise source. This assumption is often made 
    for the purpose of simplifying computation.}
    \begin{equation*}
        \noise(\vec{r}) \sim \distribution{P}_{\vec{r}},
    \end{equation*}
    where we often assume that the distribution is location independent (i.e. \(\distribution{P}_{\vec{r}} \equiv
    \distribution{P}\)).
An important property to consider when characterising noise from different sources is \emph{signal independence}, i.e. 
    whether the strength of the noise is independent from the intensity of the signal in the sense of a random variable
\begin{equation*}
    \left\{ \noise(\vec{r}) \mid \noisefree(\vec{r})  \right\} = \noise(\vec{r}).
\end{equation*}

There are many sources of noise that should be considered in the acquisition model.
The most prominent of these sources are caused by the image sensor (the \ac{CMOS}), especially in low \ac{SNR}
    acquisitions.. %  and the fluorescent proteins (\ac{EGFP}).
It is also common to regard the variations in the ambient (background) illumination i.e. the light that is not part of 
    the source we are measuring as noise.
The sources, that are the greatest contributors to the noise corruption of the expected signal are:
\marginnote{
    A popular measure used in signal processing to quantify the relative power of noise is the \ac{SNR}. For an image it is 
        defined as
        \begin{equation*}
            \SNR=\frac{\sum\limits_{u, v}\ab|\noisefree(u,v)|^2}{\sum\limits_{u, v}\ab|\noise(u,v)|^2}.
        \end{equation*}
}
\begin{marginfigure}
    \marginpdftex{poisson-distribution}
    \caption{The Poisson distribution and its continuous equivalent 
        \begin{equation*}
            \tilde{\distribution{P}}_{\lambda}(x) = \frac{e^{-\lambda}(\lambda)^x}{\Gamma(x + 1)},
        \end{equation*}
        for \(\lambda = 4\).
    }
    \labfig{pois-dist}
\end{marginfigure}
\vspace{-0.8em}
\begin{description}
    % TODO: Add figure of the relation of the SNR ratio against the signal strength in this source of noise <24-11-23> 
    % TODO: Add figure with the Poisson distribution <25-11-23> 
    \item[Photon noise]
        Inherent to the quantized nature of energy transmission of light. 
        Sensors measure irradiance by counting the number of discrete photons that land on the sensor in a given time
            interval \(\Delta \tau\). 
        The arrivals of individual photons are independent and can therefore be regarded as a Poisson process.
        The number \(N\) of photons detected in the time interval \(\Delta\tau\) can thus be described by the
            distribution~(\reffigshort{pois-dist})
            \begin{equation*}
                \distribution{P}\left[ N = x \right] = \frac{e^{-\lambda'\Delta\tau}(\lambda'\Delta\tau)^x}{x!},
            \end{equation*}
            where the rate parameter \(\lambda'\) corresponds to the expected incident photon count per unit 
            time, i.e. the expected number of photon detections per a unit of time.
        It is dependent on the source luminance.
        Since \(\Expected N = \Var(N) = \lambda'\Delta\tau \eqqcolon \lambda\), the ratio of signal to noise increases 
            with \(\sqrt{\lambda'\Delta\tau}\) and so the corruption due to this source of noise is greatest in 
            situations with low signal levels.
        For larger counts of photons, we can use the central limit theorem to approximate the photon noise by a Gaussian
            additive noise \sidecite{hasinoff2021}, whose variance and mean both depend on the signal luminance level
            \begin{equation*}
                \noise(\vec{r}) \sim \distribution{N}\left(\lambda'(\vec{r})\cdot\Delta\tau, 
                                                         \lambda'(\vec{r})\cdot\Delta\tau \right).
            \end{equation*}
    \item[Dark current]
        Semiconductor detectors in the image sensor are susceptible to thermal agitation\sidecite{aguet2009}, which can 
            cause the photocapacitors to charge even in the absence of light (hence the name dark current) leading to 
            false signals.
        Due to the phenomenon of dark current, the sensors are often cooled during acquisition to reduce this effect.
        Noise from this source is signal independent and additive.
    \item[Background noise]
        Additive and independent noise from the ambient light that reaches the sensor or unwanted light caused by other
            effects.
        An example can be the lagged emission of the fluorescent proteins.
        This effect of ambient light from outside sources is most prominent for low energy light (i.e. light with low 
            frequencies, e.g. infrared light), which is more prone to scattering and harder to effectively suppress.
        In the case of biological sensing however, a majority of what we consider background noise comes from the light 
            that is far from the focal plane and defocus causes it to blur and appear as ambient light with locally 
            (meaning in a small bounded region) limited variation.
        This is the principal reason for the use of a \ac{TIRFM} setup.
    \item[Auxiliary sources]
        A variety of causes, mainly in the imaging sensor, that are produced at the read-out time\sidenote[][*-4]{%
        The phases after the photo activation of the sensor elements (consisting of amplification, digitization of the 
            analog signal etc.).%
        }.
        Notably the charge transfer in the \ac{CMOS} sensor can be the cause of noise although, modern \ac{CMOS} sensors
            are seeing much improvement in this direction in the recent past \sidecite[*-3]{holst2014}.
        Auxiliary sources can be both signal dependent and independent.
\end{description}

If we know or assume a distribution of additive noise, it is possible to estimate its parameters using regions of the
    image with approximately or theoretically constant intensity (e.g. background regions)
    \begin{equation*}
        A \coloneq \left\{ \vec{r} \mid \noisefree(\vec{r}) \approx c  \right\}.
    \end{equation*}
\marginnote{% [*-4]
    It is often the case that we assume, or in approximation assume, that the distribution of the additive noise is 
        Gaussian and we can estimate the parameters of the noise using the for example \ac{MLE}
        \begin{gather*}
            \hat{\mu}_{\noise} = \frac{1}{|A|} \sum_{\vec{r} \in A} (\sensed(\vec{r}) - c) \\
            \hat{\sigma^2}_{\noise} = \frac{1}{|A|} \sum_{\vec{r} \in  A} (\sensed(\vec{r}) - \hat{\mu}_{\noise})^2.
        \end{gather*}
    }
Moreover, if the limitation to the theorized noise distribution is to be avoided, it can be eyeballed from the histogram
    or the histogram, kernel or other empirical estimate of the intensity distribution within \(A\) can be used.

\begin{marginfigure}
    \marginpng{otsu-threshold-leakage}
    \caption{
        The windowed Otsu threshold leads to lower intensities of some beads to be passed as background. 
        On the left, side, Otsu threshold \(t_o(\vec{r})\) is used for classifying the image pixels into background and 
            foreground by selecting the locations which have a higher intensity than \(t_o(\vec{r})\).
        On the right, the threshold is lowered to \(\tilde{t}_o(\vec{r}) \coloneqq 0.3 \cdot t_o(\vec{r})\) and the same 
            classification is performed.
        At the bottom, difference between rejected pixels by \(t_o\) and accepted by \(\tilde{t}_o\) as foreground are 
            shown.
    }
    \labfig{beads-leakage}
\end{marginfigure}

In our case, the premise that allows us to select such a set \(A\) is the lack of other non-ambient signal sources 
    besides the beads.
The problem therefore reduces to the distinction of the beads from the background, for which we can use the Otsu 
    threshold\sidecite{otsu1979}, which minimizes the intra-class variance for two classes of within the image
    \begin{multline*}
        t_o = \argmin_{t}\left( \ab|\{\sensed(\vec{r}) \mid \sensed(\vec{r}) > t\}|\cdot\Var\{\sensed(\vec{r}) \mid \sensed(\vec{r})
        > t\}\right. \\
        \left. + \ab|\{\sensed(\vec{r}) \mid \sensed(\vec{r}) \leq t\}| \cdot \Var\{\sensed(\vec{r}) \mid \sensed(\vec{r}) \leq
        t\}\right).
    \end{multline*}
It is not effective, though, to use the global Otsu threshold, because the image does not have a uniform distribution of
    intensity due to the aberrations and vignetting and partly due to the varying concentration of beads within the
    image source.
What turns out to be a better choice is a windowed Otsu threshold determined for each pixel.
Even with this enhancement though, the threshold does not achieve to completely eliminate the lower intensities at the
    edges of beads from the background which we want to determine~(\reffigshort{beads-leakage}).
Because background covers the majority of the image region and we have a lot of data points available to estimate the
    noise distribution from, we can afford to reduce the threshold for the benefit of only accepting noise at the cost 
    of rejecting some of the background in the process.
Empirically the choice was for the final lowered windowed threshold to be
    \begin{equation*}
        \tilde{t}_o = 0.3 \cdot t_o,
    \end{equation*}
    which is used in the thresholding for estimation of \(A\), the background.

A well established model of the noise distribution that is used in low \(\SNR\) acquisitions is noise consisting of two
    components that are assumed to be independent.
A Poisson distributed signal-dependent component \(\noise_{\distribution{P}}\) that represents the photon shot noise and 
    a Gaussian signal-independent component \(\noise_{\distribution{N}}\) for the read-out noise\sidecite{foi2008}.
An \(a\) and \(b\) parametrized formulation of the model that we are to estimate can be written as
\begin{equation}
    \labeq{noise-distribution-model}
    \begin{gathered}
        \begin{aligned}
            \begin{aligned}
                \noise_{\distribution{P}}(\vec{r}) &= \frac{1}{a}\cdot \alpha(\vec{r}) \\
                \noise_{\distribution{N}} &= \beta 
            \end{aligned} \quad&\quad
            \begin{aligned}
                \alpha(\vec{r}) &\sim \distribution{P}(a \cdot \noisefree(\vec{r}))\\ 
                \beta &\sim \distribution{N}(0, b^2)
            \end{aligned}
        \end{aligned}\\
        \noise(\vec{r}) = \noise_{\distribution{P}}(\vec{r}) + \noise_{\distribution{N}}.
    \end{gathered}
\end{equation}
For the estimate of this model, it is easiest to assume some form of the noise-free signal \(\noisefree\).

We selected \(A\) as the background exactly for the reason that we may make some deductions about \(\noisefree\) within
    the region of \(A\).
Unfortunately the assumption stated as a premise for selecting \(A\) as the appropriate region was too crude. 
It is not viable to assume that \(\noisefree(\vec{r}) \equiv 0\) for \(\vec{r} \in A\) according to the premise, because
    of the ambient light background noise.
We will estimate the ambient light background \(\hat{\noisefree}\) that will play the role of \(\noisefree\) within 
    \(A\) by taking a windowed mean of the background regions~(\reffigshort{ambient-light}).%
\begin{marginfigure}
    \marginpng{ambient-light}
    \caption{The estimate \(\hat{\noisefree}\) using the windowed mean of the background.}
    \labfig{ambient-light}
\end{marginfigure}%
There is an unequivocal scheme for roughly verifying the sensibility of the 
    model~(\refeqshort{noise-distribution-model}) taking the advantage of additivity of the Poisson and Gaussian 
    distributions
    \begin{align}
        \labeq{gaussian-poisson-additivity}
        \begin{gathered}
            X \sim \distribution{N}(0, \sigma^2_X),\, Y \sim \distribution{N}(0, \sigma^2_Y)\\
            X \sim \distribution{P}(\lambda_X),\,Y \sim \distribution{P}(\lambda_Y)
        \end{gathered}
        &\begin{gathered}
           \,\implies\,\\
           \,\implies\,
        \end{gathered}
        \begin{aligned}
            X + Y &\sim \distribution{N}(0, \sigma^2_X + \sigma^2_Y) \\
            X + Y &\sim \distribution{P}(\lambda_X + \lambda_Y).
        \end{aligned}
    \end{align}
    The verification is done by inspecting the histogram of \(\sensed(\vec{r}) - \hat{\noisefree}(\vec{r})
    = \hat{\noise}(\vec{r})\) for \(\vec{r} \in A\), which according to~\refeq{gaussian-poisson-additivity} should be a 
    mixture of the Poisson and Gaussian distributions (\reffigshort{demeaned-noise-histogram}) in the form
    % FIX: This needs a better formulation <29-12-23> 
    \begin{gather*}
        \sum_{\vec{r} \in A} \noise(\vec{r}) = \sum_{\vec{r} \in A} \noise_{\distribution{P}}(\vec{r})
        + \sum_{\vec{r} \in A} \noise_{\distribution{N}} = \frac{1}{a} \cdot \sum_{\vec{r} \in A}
        \alpha(\vec{r})
    + \sum_{\vec{r} \in A} \beta \\
\sum_{\vec{r} \in A} \alpha(\vec{r}) \sim \distribution{P}\ab(a \cdot \sum_{\vec{r} \in A} \noisefree(\vec{r})) \\
\sum_{\vec{r} \in A} \beta \sim \distribution{N}\ab(0, |A| \cdot b^2) 
    \end{gather*}
\begin{marginfigure}
    \marginpdftex{noise-distribution-hist-demeaned}
    \caption{%
        The histogram of \(\sensed(\vec{r}) - \hat{\noisefree}(\vec{r})\) for \(\vec{r} \in A\).
        The Poisson component of the histogram is apparent which should in part imply that the Gaussian component does 
            not play a large role in our acquisition (i.e. that \(b \ll a\)).%
        }
    \labfig{demeaned-noise-histogram}
\end{marginfigure}%

An estimate is made using the unbiased cumulant estimates from the data\sidecite{bahler2022}.

% TODO: Add SNR estimation and reason, why the noise is reduced if we add multiple acquisitions together <25-11-23> 

It is practical to assume that the, yet unknown, source signal \(\source(\vec{r})\), is directly proportional to the
    concentration of the fluorescent molecules, which in turn is directly proportional to their illuminated volume.
For this, it is necessary, that, besides the background noise\sidenote{Which can be reasonably well removed, but it is 
    important to note that only the power of the noise, linked to the mean value, is effectively reduced, not the 
    variation as such. Methods of removing the background noise generally remove the low spatial frequencies, however 
    the background noise is also prone to the photon shot phenomenon and other signal dependent noise sources, which 
    reside even after the background removal.}, to disregard the noise sources, which arise at the source 
    level.
This is a reasonable proposition, owing to the fact, that most of the noise can be attributed to the sensor and is added
    to the sensed signal after the pass through the optical apparatus.

\subsection{The Transfer Function}%
\labsubsec{the-transfer-function}

Next step is to consider how the source signal \(\source(\vec{r})\) is affected by the pass through the optical system.
This has been developed in~\refch{transfer-functions}.
As a reminder, in general, for a linear system \(\LinearSystem\), the Fredholm equation of the first kind applies \(\forall \vec{r}
    \in \R^2\)
    \begin{multline} \labeq{non-aplanatic-psf-light-formation}
        \noisefree(\vec{r}) = \LinearSystem\ab[\source(\vec{r})] = \LinearSystem\ab[\iint_{\R^2} \source(\vec{x}) \dirac(\vec{r} - \vec{x})
        \odif{\vec{x}}]  \\
            = \iint_{\R^2} \source(\vec{x}) \LinearSystem\ab[\dirac(\vec{r} - \vec{x})] \odif{\vec{x}} = \iint_{\R^2} \source(\vec{x}) \psf(\vec{r} -\vec{x}; \vec{r}) \odif{\vec{x}},
    \end{multline}
    or without the focal plane simplification \(\forall\vec{r} \in \R^3\)
    \begin{equation*}
        \noisefree(\vec{r}) = \iiint_{\R^3} \source(\vec{x}) h_3(\vec{r} - \vec{x}; \vec{r}) \odif{\vec{x}},
    \end{equation*}
    where we are using the same notation as in \refeq{additive-noise-model} and
    \begin{equation*} \labeq{PSF-in-fredholm}
        \psf(\vec{r} - \vec{x}; \vec{r}) \coloneq \LinearSystem\ab[\dirac(\vec{r} - \vec{x})]
    \end{equation*}
    denotes the \ac{PSF}.

Estimating the \ac{PSF} \(h\) that is location dependent like in \refeq{non-aplanatic-psf-light-formation} is badly
    conditioned, so we need to restraint ourselves to a model or a simplification of the transfer.
One option to handle this, is to make the assumption that the system is aplanatic%
\sidenote{%
    An optical system, or more specifically the lens configuration is aplanatic, if it is free of both spherical and 
        coma aberrations~(\refchshort{transfer-functions}).
    The aplanatic system assumption is not fully justified and there is a slight location dependence in practice.
    This issue can be addressed by estimating the \ac{PSF} on patches of the sensed image separately. 
    This will be done in later sections.
    % TODO: Closer reference <24-11-23> 
}.
This assumption makes the \ac{PSF} focal plane location independent i.e. \(\forall\vec{x},\vec{r} \in \R^2\)
    \begin{equation*} \labeq{aplanar-psf-light-formation}
        \psf(\vec{r} - \vec{x}; \vec{r}) = \psf(\vec{r} - \vec{x})
    \end{equation*}
     within the same \(z\)-offset plane from the focal plane%
\sidenote{%
    However it is important to realize that it is not true that under the assumption of an aplanatic system that 
        \ac{PSF} \(h_3\) is identical as a function \(\R^3 \to \R\) for the whole imaged volume.
    The lack of spherical aberration and coma only fixes \(h\) on the \(x\) and \(y\) axes of the object space, so for 
    \(\vec{x},\vec{r} \in \R^3\)
    \begin{equation*}
       h_3(\vec{r} - \vec{x}; \vec{r})  \equiv h_3(\vec{r} - \vec{x}; z).
    \end{equation*}
}.

For the location invariant \(h\), we can rewrite \refeq{non-aplanatic-psf-light-formation} as
    \begin{equation}
        \labeq{convolution-psf}
        \noisefree(\vec{r}) = \iint_{R^2} \source(\vec{x}) \psf(\vec{r} - \vec{x}) \odif{\vec{x}} = 
            (\source \conv \psf)(\vec{r}),
    \end{equation}
    where \(\conv\) denotes the convolution operation.
The locality\sidenote{Locality is understood here as the property that if \(\supp(\psf) \subset 
    \ball_\alpha\), then for a source \(\source\) for which \(\supp(\source) \subset \ball_\beta\) we
    there exists a \(\gamma\) such that \(\supp(\source \conv \psf) \subset \ball_\gamma.\)} of the convolution is key to making the 
    \ac{PSF} more intuitive and easier in practice to measure.
It is, because unlike the \ac{OTF}, we are not forced to make the estimate from the full image or strenuously map the
    frequency domain of the parts used for measurement onto the frequencies of the acquisition image.
% TODO:  What changes when the PSF changes...  <22-11-23, kunzaatko> 
% TODO: What changes when the OTF changes <22-11-23> 
Consequently, parts that are easiest to work with can be selected.
% TODO: Another condition for this to work is that the PSF has a bounded support with a bound that is limited by the
% region sizes that we select. This should be argued. It is not the case but the PSF has is small enough to consider
% this as true <27-11-23> 
This can be beneficial in eliminating noisy or inconsistent data, or data, that our model does not consider.


\subsection{The Source}
\labsubsec{the-source}
In order to devise \(\psf\) from \refeq{convolution-psf}, it is necessary to determine \(\source(\vec{r})\), for the
    parts that are used for the measurement.
Consider what it would would mean for the acquisition, if \(\source= \sum_{\vec{a} \in A} \dirac_{\vec{a}}\).
The sensed image, where noise is disregarded, would be
    \begin{equation}
        \labeq{noisefree-with-sum-of-dirac-sources-psf-general}
        \noisefree(\vec{r}) = \ab(\source \conv \psf)(\vec{r}) = \ab(\ab(\sum_{\vec{a} \in A} \dirac_{\vec{a}}) \conv
            \psf)(\vec{r}) = \sum_{\vec{a} \in A} \psf(\vec{r} - \vec{a}),
    \end{equation}
    and if adapted to the more natural notation, where we emphasize, which point underlines the term of the expansion in
    the image formation formula, i.e. \(\psf(\vec{r} - \vec{a}) \eqqcolon \psf_{\vec{a}}
    (\vec{r})\), \refeq{noisefree-with-sum-of-dirac-sources-psf-general} changes to
    \begin{equation}
        \labeq{noisefree-with-sum-of-dirac-sources-psf-natural}
        \noisefree(\vec{r}) = \sum_{\vec{a} \in A} \psf_{\vec{a}}(\vec{r}),
    \end{equation}
    and we see that the noise-free sensed signal is merely a set of copies of the \ac{PSF} at the locations of the
    initial \(\dirac\) sources.
\marginnote[*-5]{\refeq{noisefree-with-sum-of-dirac-sources-psf-natural} is also applicable for the non-aplanatic 
    transfer, but it is important to realize that in that case the index \(\vec{a}\) of the \ac{PSF} does not denote 
    only a change in coordinates (translation by \(\vec{a}\)), but also the change in the function \(\psf_{\vec{a}}\) 
    itself.} %  This can be used alongside partitioning the sensed image into regions and measuring the relation of the 
    % \ac{PSF} to the location with in the \refsec{location-dependent-psf}}

In practice, this type of source is generated by imaging microspheres most commonly made from plastics, stained with 
    fluorescent dyes.
Acquisitions of fluorescent beads is often used by biologists for many purposes such as calibrating equipment, because 
    it is possible to manufacture them with a high degree of precision and modeling the source that is generated by them
    is straightforward as we will observe in~\refsec{source-model}.

An important aspect that is essential to be able to measure the \ac{PSF} in a direct manner is that the diameter of the
    bead is sub-diffraction-resolution\sidenote[][*-3]{It is possible, and even beneficial if the diameter is not
    sub-pixel-resolution, but it does not make a difference in the principle of the \ac{PSF} measurement.}.
If the diameter were not to be sub-diffraction-resolution the source signal would be different from the
    \(\dirac_{\vec{a}}\) of \refeq{noisefree-with-sum-of-dirac-sources-psf-general}, but rather some other support 
    limited source for which a model \(\hat{\source}_{\vec{a}}\) would be necessary to develop, and the noise-free
    sensed signal \(\noisefree\) would not be direct copies of the \ac{PSF} \(\psf\) but
    \begin{equation}
        \labeq{non-sub-diffraction-resolution-source}
        \noisefree(\vec{r}) = \sum_{\vec{a} \in A} \ab(\hat{\source}_{\vec{a}} \conv \psf)(\vec{r})
    \end{equation}
    instead.
Deriving the \ac{PSF} \(\psf\) from the signal in \refeq{non-sub-diffraction-resolution-source} would require 
    deconvolution of the individual components \((\hat{\source}_{\vec{a}} \conv \psf)\) with the assumed model 
    \(\hat{\source}_{\vec{a}}\) and due to noise, this would necessarily mean that we would need to make further 
    assumptions either about the deconvolved \(\hat{\psf}\) or the noise.

The acquisition that is used for the \ac{PSF} estimation is done by imaging a suspension of \emph{TetraSpeck\TM 
    microspheres, \(\qty{0.1}{\micro\metre}\)}, where the diameter is 
    \(\qty{0.099 \pm 0.008}{\micro\metre}\)\sidecite{smith2020}, which is safely below the diffraction resolution of 
    \begin{equation*}
        \frac{\lambda_{\text{em}}}{2\NA} = \frac{\qty{488}{\nano\meter}}{2\cdot 1.4} \approx \qty{174}{\nano\meter}
        = \qty{0.174}{\micro\metre}.
    \end{equation*}

% TODO: Figure with selected beads marked from the acquisition <28-11-23> 

% TODO: Either in this section or in the section about the Abbe resolution limit, there should be the two point
% resolution definition <27-11-23> 


\section{\ac{PSF} from  Beads}%
\labsec{psf-from-beads}

It is common in among biologists to select the sensed signal produced by one bead near the centre of the acquisition 
    image and pronounce it as the \ac{PSF} of the optical system.
% TODO: References (these references should be those which are used elsewhere 
% (perhaps in the SIM reconstruction papers)) <28-11-23> 
There is ample room for improvement from this approach. 
Noteworthy drawbacks of this approach are that
    \begin{enumerate}[label=(\alph*)]
        \item noise contained in a single bead sample can be substantial, \labitm{noise}
        \item the bead centre may not be exactly located at the centre of a pixel and therefore the \ac{PSF} might be 
            biased, \labitm{pixel-centring}
        \item it disregards the dependence on the \(z\)-offset from the focal plane of the particular bead 
            sample used, \labitm{z-offset}
        \item it is negligent towards the potential dependence of the sampled bead position within the focal plane.
            \labitm{position-dependence}
    \end{enumerate}

A simplest alternative that addresses all of the aforementioned drawbacks and fixes most of them to certain extent is 
    averaging multiple bead samples.
The steps of performing an estimate from multiple bead samples are:
\begin{enumerate}[label=\roman*.]
    % TODO: Add a figure of the beads and the focal plane <28-11-23> 
    % TODO: Add a mathematical explanation based on defocus "operator" (diffusion), why the same model of beads
    % leads to a lower centre intensity if it is shifted by some Δz <28-11-23> 
    \item Selecting approximate centres \(C\) of isolated beads with the highest intensity%
        \sidenote[][*-3]{%
            Beads with the highest intensity can be assumed to have the tip closest to the focal plane of the optical 
                system.
            Since we are looking for the \ac{PSF} at the focal plane (i.e. \(z = 0\)), these beads are selected for the 
                estimate.
        }.
            % \[
            %     C \coloneqq \ab\{\vec{c}_i \coloneqq (u_i, v_i) \mid i \in \hat{n}\}.
            % \] 
    \item Cutting a patch around each centre \(\vec{c}_i\) containing the whole bead sample.
        % \begin{gather*}
        %     U_i \coloneqq \ab\{u_i - \Delta u, \ldots, u_i + \Delta u\} \\
        %     V_i \coloneqq \ab\{v_i - \Delta v, \ldots, v_i + \Delta v\} \\
        %     P \coloneqq \ab\{ P_i \coloneqq (\sensed_{uv})_{u \in U_i, v \in V_i} \mid i \in \hat{n} \}.
        % \end{gather*}
    \item Subtracting the ambient light background from each patch. 
        The intensity of the ambient light in the location of the bead is estimated by the \(0.15\) quantile of the 
            intensities of a larger patch around the bead sample.
    \item Finding the sub-pixel centres of the beads by \ac{LSE} fitting a Gaussian%
        \sidenote[][*-5]{%
            Any other unimodal function, which has its mode at the origin would suffice.
            The Gaussian is strategic here, since the \ac{PSF} has close shape to the Gaussian.
            In fact, it is often modelled by a it.
        }
        to the sample array \labitm{centring}
    \item Shifting by sub-pixel offsets of the estimated beads centres from \refitm{centring} to fix the them to a whole
        pixel value via Fourier shift theorem.
    \item Averaging the patch pixels to obtain the \ac{PSF} estimation.
\end{enumerate}

We subtract the backgrounds of the sample beads before averaging them, therefore, we may assume that their noise is 
    centred, i.e. \(\Expected\noise =~0\).
The variance of the \ac{PSF} values contained in the estimate is reduced regardless of the distribution of the noise 
    in the sensed samples due to the law of large numbers.

We will perform the estimate of the hypothesised focal plane global \ac{PSF} on bead samples from the center of the
    image~(\reffigshort{selected-beads-CC}).
The cut-outs along with the fitted centres of the bead samples are shown in~\reffig{selected-beads-patches-CC}.
The final \ac{PSF} estimate and its comparison with the model introduced in~\refch{sim-reconstruction} is visible 
    in~\reffig{psf-avg-estimate-CC}.

\begin{marginfigure}[*-15]
  \marginpng{selected-beads-CC}
  \caption{Bead samples selected from the center of the image. For better accuracy the locations were selected from the
  \ac{HR} reconstructed image.}
  \labfig{selected-beads-CC}
\end{marginfigure}

\floatsetup[figure]{ % Captions for figueres
    capbesideposition={top,outside},%
}%
\begin{figure}[h]
    % \png{selected-beads-patches-CC-noticks}
    \png{selected-beads-patches-CC}
    \caption{The bead samples cut-outs from the center of the image and their fitted centres.}
    \labfig{selected-beads-patches-CC}
\end{figure}
\floatsetup[figure]{ % Captions for figueres
    capbesideposition={bottom,outside},%
}%

% TODO: Here figure with one cutout before and after centering <30-12-23> 

\begin{marginfigure}[*-7]
    \marginpng{model-vs-avg-psf}
    \caption{%
        The \ac{PSF} estimate versus the model. 
        The model is clearly wider, which implies that the diffraction-resolution was overestimated by the model which
            is in line with the reconstruction flaws we discussed in \refsec{deconvolution-reconstruction}.
    }
    \labfig{psf-avg-estimate-CC}
\end{marginfigure}

A first test to see if the estimate is better than the proposed model is to compare the deconvolution of the \ac{LR}
    image with the proposed \ac{PSF}~(\reffigshort{model-vs-avg-deconv-CC}).
The deconvolution in this case is performed by formulating an optimization problem very similar to the formulation of 
    the optimization problem of reconstruction~(\refsecshort{inversion-reconstruction}), with the sole
    difference that the forward model is the convolution from~\refeq{convolution-psf} and no regularization is 
    performed.
Importantly for assessing the accuracy of the deconvolution, we need to realize, what an ideal deconvolution should look
    like.
The key is that the beads are sub-diffraction-resolution and as such, they should have a single intensity and should be
    comprehended by the optical system as a single point as defined by the two-point resolution, but they do not have a 
    sub-pixel-resolution so they are supposed to be visible over multiple pixels.
The sampling rate of our setup is \(\qty{61}{\nano\meter}\) so, since the beads have a diameter of
    \(\approx\qty{99}{\nano\meter}\), they should have a image plane diameter of about \(\qty{1.6}{\pixel\lowres}\).

\pagebreak
\begin{marginfigure}[*+15]
    \marginpng{model-vs-avg-psf}
    \caption{%
        The \ac{PSF} estimate versus the model. 
        The model is clearly wider, which implies that the diffraction-resolution was overestimated by the model which
            is in line with the reconstruction flaws we discussed in \refsec{wiener-filter-reconstruction}.
    }
    \labfig{psf-avg-estimate-CC}
\end{marginfigure}


\begin{figure}[h]
    \png{model-vs-avg-psf-deconv-CC-vs-sensed}
    \caption{%
        Deconvolution of the \ac{LR} images performed with the \ac{PSF} model and the \ac{PSF} estimate.
        The beads size (\(\qty{99}{\nano\meter}\)) in the \ac{LR} sampled (\(\qty{61}{\nano\meter/\pixel}\)) source 
            should be approximately \(\qty{1.6}{\pixel\lowres}\).
        The possible explanation for the deviation from this expectation is the varying \(z\)-offset from the focal plane 
            which of the individual beads that causes defocus.
    }
    \labfig{model-vs-avg-deconv-CC}
\end{figure}

% TODO: Here fig of the inds where we estimate the location dependent PSF with the beads selected in them <30-12-23> 
% TODO: Here comparison of the psf estimates from the sections of the image <30-12-23> 

% TODO: Here for the estimate of the noise variance of the estimate from the determined distribution of noise <29-12-23> 
% TODO: Here fig of the selected beads for all patches <29-12-23> 

% SSIM s modelem
% TODO: Here comparison of the parameters of the fit from different locations of the object field <31-12-23> 

\section{The Source Model}
\labsec{source-model}

For establishing a model of the source signal, we need to
\begin{enumerate}[label=\roman*.]
    \item \labitm{bead-model} determine the model of a single bead sampled in the destinations image sampling rate (i.e. for
        \ac{LR} images the sampling rate is \(\qty{61}{\nano\meter\per\pixel\lowres}\) and for \ac{HR} images the
        sampling rate is \(\qty{30.5}{\nano\meter\per\pixel\highres}\))
    \item \labitm{locations-of-beads} find the bead locations within the object field and create the full image model as
        a sum of the single bead model from \refitm{bead-model} centred at each location.
\end{enumerate}

\subsection{Single Bead Model (\refitmshort{bead-model})}%
\labsubsec{single-bead-model}

The beads are effectively balls that are suspended at some \(z\)-offset from the focal plane within the object field. %
% TODO: Here a figure of the beads ball and the focal plane that touches the tip of the balls <31-12-23> 
We will model them as such balls with a uniform distribution of fluorescent particles at their surface.
A simplification that is made is the, not fully accurate, assumption that the \(z\)-offset of all the beads is
    \(\qty{0}{\nano\meter}\).
One consideration must be kept in mind and that is the \ac{TIR} nature of the acquisition.
This means that the illumination intensity, more befittingly termed excitation intensity in this context, 
    \(\ip_{\text{ex}}\) is according to the theory from \refch{light} attenuated by \(e^{\alpha z}\) where \(\alpha\) is 
    the attenuation constant that is determined by the super-critical angle of the illumination and \(z\) denotes the 
    \(z\)-offset from the focal plane of the location of the point in the object field. %
% TODO: Closer reference <31-12-23> 
A full formulation of \(\ip_{\text{ex}}\) is therefore
    \begin{equation*}
        \ip_{\text{ex}}(x,y,z) = \ip_{\text{ex}}(x,y,0) \cdot e^{\alpha z},
    \end{equation*}
    where \(\ip_{\text{ex}(x,y,0)}\) is constant in our case and can be set to equal \(1\) without loss of generality to
    the bead model.

The \(z\)-offset of a given patch of the surface of a ball with its tip at the point \((0,0,0)\) is 
    \begin{equation*}
      z(x,y) = R - \sqrt{R^2 - (x^2 + y^2)},
    \end{equation*}
    where \(R\) is the radius of the ball.

We are working with sub-saturation excitation intensities and time intervals, which means that
    \begin{equation*}
        \ip_{\text{ex}}  \cdot d_{\text{fl}} \propto \ip_{\text{em}},
    \end{equation*}
    where \(\ip_{\text{em}}\) is the emission intensity and \(d_{\text{fl}}\) is the distribution of the
    fluorescent particles in the object field which was already assumed to be constant for the beads surfaces.
This means that the emission intensity \(\ip_{\text{em}}\) of a bead at the point \((x,y)\) in the focal 
    plane~(\reffigshort{bead-model}) is
    \begin{equation*}
        \ip_{\text{em}}(x,y) = \ip_{\text{ex}}(x,y,z) = \ip_{\text{ex}}(x,y,0) \cdot e^{\alpha z(x,y)}
            = \frac{e^{\alpha R}}{ e^{\alpha\sqrt{R^2 - (x^2 + y^2)}}}.
    \end{equation*}
It is crucial to emphasize another simplification that is made and is quite comical considering the goal of developing
    this model.
We are projecting the emission intensity from the \((x,y,z)\) location directly to the emission intensity at the 
    \((x,y)\) location on the focal plane.
This is not correct, due to the \(z\)-offset plane not having the same transfer function as the focal plane.
The major impact of this omission is that we do not consider defocus (as part of the \ac{PSF} difference) in any way.

The sampling in done by integration of the point emission function \(\ip_{\text{em}}(x,y)\) over the sensor element that
    corresponds to a given pixel
    \begin{equation*}
        m(u,v) \coloneqq \ip_{\text{em}}(u,v) = \iint\limits_{\mathclap{u - \tfrac{\Delta_{x}}{2},\, v - \tfrac{\Delta_{y}}{2}}}^{\mathclap{u + \tfrac{\Delta_{x}}{2},\, v + \tfrac{\Delta_{y}}{2}}} \ip_{em}(x,y) \odif{x} \odif{y},
    \end{equation*}
    where \(\Delta_{x}\) and \(\Delta_{y}\) are the dimensions of the sensor element (i.e. the sampling rates in the
    \(x\) and \(y\) axes).

\begin{marginfigure}
    \marginpng{source-model-single-bead}
    \caption{\ac{LR} and \ac{HR} bead models. The radius of the bead in the \ac{LR} sampling rate is \(\approx
        \qty{1.6}{\pixel\lowres}\), and for the \ac{HR} bead is is \(\approx\qty{3.2}{\pixel\highres}\).}
    \labfig{bead-model}
\end{marginfigure}
    

\section{Locating the Beads (\refitmshort{locations-of-beads})}%
\labsec{locating-the-beads}

The locations of the beads are determined by an empirically constructed algorithm that is inspired by the general Hough
    transform and image feature location algorithms.
The idea is to
    \begin{enumerate}[label=\roman*.]
        \item select locations that are candidate to the positions of beads, using a global algorithm with a very low 
            threshold such that no true location of a bead gets discarded, \labitm{global-candidates}
        \item filter the candidate locations by a local algorithm to discard the false positives \labitm{local-filter}
        \item determine the final sub-pixel locations of the beads \labitm{final-location}.
    \end{enumerate}
The reason for finding the initial candidates by a global algorithm~(\refitmshort{global-candidates}) is induced by the 
    preference of an algorithm that has a simple formulation and for which there is not a overwhelming number of 
    empirical parameters that have to be set.
The local filtering of the locations~(\refitmshort{local-filter}) is motivated by the fact that the bead samples are 
    typically locally dependent and therefore to robustly discard the false positive locations cannot to be done by 
    comparing all the locations, but only the ones that are comparable with their properties to the one that is being 
    judged, i.e. the locations in the neighbourhood of the bead in consideration.


% \begin{marginfigure}
%     \marginpng{source-model-ccorr}
%     \caption{%
%         Normalized cross-correlation of the Weiner filter reconstructed \ac{HR} image with the \ac{HR} bead model 
%             \(m_{\ac{HR}}\).
%         }
%     \labfig{source-model-ccorr}
% \end{marginfigure}

\floatsetup[figure]{ % Captions for figueres
    capbesideposition={top,outside},%
}%
\begin{figure}
    \png{source-model-ccorr}
    \caption{%
        Normalized cross-correlation \(\correlation_{m_{\ac{HR}} \times \reconstructed}\) of the Weiner filter 
        reconstructed \ac{HR} image with the \ac{HR} bead model \(m_{\ac{HR}}\).
        }
    \labfig{source-model-ccorr}
\end{figure}
\floatsetup[figure]{ % Captions for figueres
    capbesideposition={bottom,outside},%
}%

The global algorithm that is used for selecting the candidates is based on thresholding the cross-correlation of the
    \ac{HR} reconstruction \(\reconstructed\) with the \ac{HR} bead model 
    \(m_{\ac{HR}}\)~(\reffigshort{source-model-ccorr}) in the form
    \begin{equation*}
        \correlation_{m_{\ac{HR}} \times \reconstructed}(u,v) = 
        \frac{\sum\limits_{m,n}\reconstructed(u + m,v + n) \cdot m_{\ac{HR}}(m,n)}{\sqrt{\sum\limits_{m,n} m^2_{\ac{HR}}(m,n)
        \sum\limits_{m,n}\reconstructed^2(u + m,v + n)}},
    \end{equation*}
    where the limits of summation of \(m,n\) are taken to be the overlap regions of the \((m,n)\)-shifted 
    \(\reconstructed\) and \(m_{\ac{HR}}\) and \(\reconstructed\) is constantly replicated beyond its sampled borders by 
    the border value\sidecite[*-23]{gonzalez2008a} in order to account for \((u,v)\) in the full sampled
    area\sidenote[][*-21]{%
    It is often the case that we instead perform the correlation of the demeaned (centred) functions \(f - \bar{f}\)
    and \(g - \bar{g}\) in which case the correlation has the meaning of the correlation coefficient and the has values
    inclusively between \(-1\) and \(1\).}.
The areas where the cross-correlation \(\correlation_{m_{\ac{HR}} \times \reconstructed}\) is above \(\alpha = 0.62\) 
are considered to be candidates for beads locations (\reffigshort{ccorr-discarded})
    \begin{equation*}
        \correlation_{m_{\ac{HR}} \times \reconstructed}(u,v) > \alpha.
    \end{equation*} % TODO: Here margin figure Correlation above <31-12-23> 


\begin{marginfigure}[*-20]
    \marginpng{source-model-ccorr-discarded}
    \caption{%
        The cross-correlation to establish the beads locations candidates.
        It is visible that some locations within the background are passed as candidates. 
        This is even more pronounced in the distal areas of the image near the edges.
    }
    \labfig{ccorr-discarded}
\end{marginfigure}


The filtering of the candidate locations (\refitmshort{local-filter}) is then made by comparing the locations to a wide
and narrow neighbourhood mean intensity of the image, \(\reconstructed_{\text{W}}\) and
    \(\reconstructed_{\text{n}}\) respectively.
In order for the bead \ac{ROI} to be considered as a true bead \ac{ROI} and pass to the next filtering step, the ratio
    of its \ac{ROI} intensities to the mean intensities must be above some threshold.
The comparison with the wide mean of intensity \(\reconstructed_{\text{W}}\) is taken with the motivation of 
    discarding the false positive locations that are due to noise that is reminiscent of the bead model.
It is made by a windowed mean with a window size of \(201 \times 201~\unit{\pixel\highres}\) and the ideal threshold is empirically determined 
    to be \(\beta_\text{W} = 8.0\). 
The final filter (\reffigshort{wide-discarded}) of the candidate \acp{ROI} by the wide mean is done as
\begin{equation*}
  \reconstructed(u,v) > \beta_\text{W} \cdot \reconstructed_{\text{W}}(u,v).
\end{equation*}

\begin{marginfigure}[*-15]
    \marginpng{source-model-wide-mean-discarded}
    \caption{%
        The wide mean filter of the candidate bead \acp{ROI}.
        It is effective in discarding the false positive locations that are due to the background noise and therefore
            have a low intensity in comparison with the mean of their neighbourhood intensities.
    }
    \labfig{wide-discarded}
\end{marginfigure}

The comparison with the narrow mean \(\reconstructed_{\text{n}}\) is performed in order to discard the false 
    positive locations that are present because of the reconstruction artefacts that occur using the Wiener filter
    reconstruction~(\refsecshort{wiener-filter-reconstruction}) and to discard the pixels connecting \acp{ROI} of beads
    that are close together, which is necessary for the \refitm{final-location}. % TODO: Closer reference <04-01-24> 
It is made by a windowed mean with the window size of \(15\times 15~\unit{\pixel\highres}\) and the empirical threshold 
    is determined to be \(\beta_\text{n} = 2.46\)
Similarly to the wide mean threshold, the filtering of the \acp{ROI} is done as
    \begin{equation*}
    \reconstructed(u,v) > \beta_\text{n} \cdot \reconstructed_{\text{n}}(u,v).
    \end{equation*}
The final narrow mean filter that is used on the \acp{ROI} of the candidate beads is in \reffig{narrow-discarded}
\begin{marginfigure}[*-13]
    \marginpng{source-model-narrow-mean-discarded}
    \caption{%
    In the wide mean, the \acp{ROI} of these beads remain connected.
    This is not desired for in the next step, where the centroids of the \acp{ROI} are used as initial locations of the 
        candidate beads for the Gaussian fit.
    A filter that has a narrow window size is efficient in reducing these connections.
    Another possibility would be to use the morphological operation of \emph{thinning}\sidecite[*+3]{gonzalez2008b} on the binary image
        after the wide mean filter to remove the connections.
}
    \labfig{narrow-discarded}
\end{marginfigure}
% TODO: Here margin figure. The discarded locations by narrow mean <04-01-24> 

The sub-pixel center localization~(\refitmshort{final-location}) is performed similarly as in~\refsec{psf-from-beads} by
    fitting a Gaussian to the centroids of the individual \acp{ROI}.
It assumes that the \acp{ROI} of two neighbouring beads does not overlap, which is ensured (at a best effort bases) by 
    setting the thresholds \(\alpha\), \(\beta_\text{W}\) and \(\beta_\text{n}\) in the previous steps.
The fitted Gaussian are further utilized to filter the artefacts that are passed as beads from the previous steps.
This is done by setting a threshold on the ratio of means \(\mu\) and deviations \(\sigma\) of the fitted Gaussian in
    a neighbourhood and also by noticing that the artefacts are at set relative locations from the bead that causes
    them~(\reffigshort{weiner-reconstruction-artifacts}).

% \begin{marginfigure}
%     \marginpng{source-model-gaussian-sum}
%     \caption{}
%     \labfig{gaussian-sum-model}
% \end{marginfigure}

\begin{figure}[h]
    \png{source-model-gaussian-sum}
    \caption{%
        Comparison of the reconstructed \(\reconstructed\) along with the sum of the fit Gaussian functions for every
            candidate location and the assumed artefacts are filtered by their location and the relative fitted mean
            \(\mu\) and deviation \(\sigma\).
        Even though many of the artefact candidate locations are filtered out, the outcome is not ideal and it could be
            a better strategy to use both the \ac{HR} reconstructed image for its better resolution and the \ac{LR}
            deconvolved image, perhaps in the form of a cross-correlation with the \ac{LR} bead model to filter out the
            artefact candidates.
    }
    \labfig{gaussian-sum-model}
\end{figure}

The final source model (\reffigshort{source-model}) \(\hat{m}\) is then obtained by summing the single bead models 
    \(m_\text{\ac{HR}}\) at the filtered locations, weighted in order for the maximum intensity of the bead at the 
    location to match the intensity of the fitted Gaussian.
Weighting is done to account for the different \(z\)-offsets as mentioned in \refsubsec{single-bead-model}.

% \begin{marginfigure}
%     \marginpng{source-model}
%     \caption{}
%     \labfig{source-model}
% \end{marginfigure}

\begin{figure}[h]
    \png{source-model}
    \caption{The reconstructed \ac{HR} image and the final model of the source.}
    \labfig{source-model}
\end{figure}

\section{Estimating the PSF from the Model}%
\labsec{estimating_the_psf_from_the_model}

The estimate of the degradation function, which for an invariant (aplanatic) system corresponds to the transfer 
    function, supported that we have a source model and its acquisition can be made very trivially 
    \sidecite{gonzalez2008c}. 
If we write the noisy transfer in the Fourier domain, we get
    \begin{equation}
      \Fsensed(\vec{k}) = \otf(\vec{k}) \cdot \Fsource(\vec{k}) + \Fnoise(\vec{k}),
    \end{equation}
    where \(\otf\) is the transfer what we want to estimate.
Therefore in a noisy acquisition if \(\Fnoise\) is sufficiently small relative to \(\otf \cdot \Fsource\), the \ac{SNR} 
    is sufficiently large, then
    \begin{equation*}
        \frac{\Fsensed(\vec{k})}{\Fsource(\vec{k})} \approx \otf(\vec{k}).
    \end{equation*}
Thus, if we have the model of the source \(\hat{m} \approx \source\), it is sufficient to make the Fourier transform and 
    pixel-wise division with the sensed image in the Fourier domain and we have an estimate \(\hat{\otf}\) of the 
    \ac{OTF}.

% TODO: Here fig of the estimate that is from the full source model <05-01-24> 
% TODO: Here fig of the simple estimate of the OTF from a model from multiple locations within the image <04-01-24> 

\section{\ac{PSF} from \ac{SIM}}
\labsec{psf-from-sim}

Since our terminal goal is to perform the \ac{SIM} reconstruction using the transfer function, we can use the model
    \(\hat{m}\), and a similar optimization problem as in \refsec{inversion-reconstruction} and try to determine the
    best transfer function specifically for this purpose.
The optimization problem given a model \(\hat{m}\) of the source, given as
    \begin{equation}
        \labeq{optimization-psf}
        \hat{\otf} = \argmin_{\otf} \loss[\sensed,  \fwd[\otf]],
    \end{equation}
The only change from \(\fwd\) in \refsubsec{forward-model} in \refeq{optimization-psf} is that now we view the \(\otf\) as the variable
    instead of the source \(\source\) and we substitute the source in the forward mapping \(\fwd\) by \(\hat{m}\).
We can use the same loss \(\loss\) in \refeq{optimization-psf} as in \refsubsec{loss-function}.
It is not necessary to perform regularization if instead of the \ac{OTF} \(\otf\) the optimization done for the \ac{PSF}
    \(\psf\) which can be sampled over a small region and therefore we can avoid the optimization problem to be
    underspecified, since we have many more data points within the data \(\sensed\) than we want to optimize over in the
    \ac{PSF} \(\psf\).

% TODO: Here reconstruction through optimization of the global and for the individual parts <05-01-24> 

\section{Revised Reconstruction}%
\labsec{revised-reconstruction}

\section{Discussion}%
\labsec{discussion}

% TODO: This section should probably be removed and the introduction should only state the possibility and cite Fourier
% Optics instead of implementing this. <31-12-23> 
% \section{Phase Retrieval}%
% \labsec{phase-retrieval}
